
We continue under the hypotheses of \SGARefLink{sga3:xiv:2.4}{2.4}. Let
\[
  \mathcal D\colon
  \bigl((\mathrm{Sch})_{/S}\bigr)^\circ
  \longrightarrow \mathrm{Sets}
\]
be the functor defined by
\[
  \mathcal D(S')
  =\left\{
    \text{Cartan subalgebras of }\mathfrak g_{S'}
   \right\}.
\]
Introduce also the functor
\[
  X(S')
  =\left\{
    (\mathfrak d,a)\ \middle|\
    \begin{array}{l}
      \mathfrak d\text{ is a Cartan subalgebra of }\mathfrak g_{S'},\\
      a\text{ is a section of }\mathfrak d
    \end{array}
   \right\}.
\]
The two projections
\[
  (\mathfrak d,a)\longmapsto\mathfrak d
  \qquad\text{and}\qquad
  (\mathfrak d,a)\longmapsto a
\]
give morphisms
\[
  p\colon X\longrightarrow\mathcal D
  \qquad\text{and}\qquad
  \psi\colon X\longrightarrow W(\mathfrak g).
\]

When \((\mathrm C_0)\) holds, let \(U\) also denote the open subset
of regular points of \(W(\mathfrak g)\) \textup{(\SGARefLink{sga3:xiv:2.10}{2.10})}. Condition
\((\mathrm C_2)\) is then equivalent to saying that the morphism
\[
  \psi^{-1}(U)\longrightarrow U
\]
induced by \(\psi\) is an isomorphism; a priori it is a
monomorphism by \SGARefLink{sga3:xiv:2.6}{2.6}.

It is immediate that
\[
  p\colon X\longrightarrow\mathcal D
\]
is representable by a projection of vector bundles. More
explicitly, for every \(S\)-morphism \(S'\to\mathcal D\)
corresponding to a Cartan subalgebra \(\mathfrak d\) of
\(\mathfrak g_{S'}\), the fiber product
\[
  X\times_{\mathcal D}S'
\]
is represented by the vector bundle \(W(\mathfrak d)\) over
\(S'\). Thus, if \(\mathcal D\) is representable, then so is \(X\).
We have:

\paragraph{\textbf{Theorem 2.16.}}
\SGARefTarget{sga3:xiv:2.16}
{\itshape Let \(S\) be a prescheme and let \(\mathfrak g\) be a Lie
algebra over \(S\) whose underlying \(\mathcal O_S\)-module is
locally free of finite type. Suppose that condition
\((\mathrm C_0)\) of \SGARefLink{sga3:xiv:2.9}{2.9} holds.

\smallskip
\noindent
\textup{a)}
The functor \(\mathcal D\) of Cartan subalgebras of
\(\mathfrak g\), defined above, is represented by a
quasi-projective prescheme of finite presentation over \(S\).
The same holds for the functor \(X\).

\smallskip
\noindent
\textup{b)}
If condition \((\mathrm C_2)\) of \SGARefLink{sga3:xiv:2.9}{2.9} holds, then
\(\mathcal D\) and \(X\) are smooth over \(S\), and the morphism
\[
  \psi^{-1}(U)\longrightarrow U
\]
induced by \(\psi\) is an isomorphism.

\smallskip
\noindent
\textup{c)}
Continue to suppose that \((\mathrm C_2)\) holds. Let \(s\in S\),
and let \(\mathfrak d_0\) be a Cartan subalgebra of
\(\mathfrak g(s)\), corresponding to a \(k(s)\)-rational point
\(d\) of \(\mathcal D\). Suppose that \(\mathfrak d_0\) contains a
regular element of \(\mathfrak g(s)\), a condition automatically
satisfied when \(k(s)\) is infinite. If \(r\) is the infinitesimal
rank of \(\mathfrak g(s)\) and \(n\) its rank over \(k(s)\), then
there exists an open neighborhood \(V\) of \(d\) in
\(\mathcal D\) that is \(S\)-isomorphic to an open subset \(V'\)
of
\[
  S[t_1,\ldots,t_{n-r}].
\]}
